\section{Исследовательский раздел}
В данном разделе будут описаны проведеные исследование временной эффективности
разработанной системы, исследуемыми параметрами являются время выполнения
анализа каскадных сбоев в зависимости от количества узлов системы и время
поиска циклов в системе.

\subsection{Технические характеристики}
Замеры времени проводились на ноутбуке со следующими характеристиками:
\begin{itemize}
    \item ядро операционной системы --- GNU/Linux 7.0.10\_arch1-1 x86\_64~\cite{ref:archbtw};
    \item объём оперативной памяти --- 32 ГБ;
    \item процессор --- Intel\textregistered \enspace 
    Core\texttrademark \enspace Ultra 7 155H с тактовой частотой 4.80 ГГц;
    \item количество физических ядер --- 16;
    \item количество логических ядер --- 22.
\end{itemize}

При замерах процессорного времени ноутбук был подключён к сети электропитания.
Во время тестирования процессор был загружен только выполнением системных
процессов, а также системы тестирования.

\subsection{Технология замеров времени}
Для каждого описанного далее исследования проводилось 100 замеров времени для
каждого набора входных данных, полученные результаты усреднялись. Для получения
времени выполнения запроса использовались временные метки между началом и концом
выполнения запроса. Временные метки получены из СУБД Neo4j, результирующим временем
считается разница времени между метками.

\subsection{Первое исследование}
\subsubsection{Описание исследования}
Данное исследование направлено на выявление зависимости времени выполнения поиска
каскадных зависимостей от количества узлов в системе. Условия исследования:
\begin{itemize}
    \item количество циклов в графе системы зафиксировано;
    \item количество узлов в графе системы изменяется от 1000 до 10000 с шагом 100;
    \item количество связей в графе системы равняется количеству узлов, умноженному на два, на каждом шаге.
\end{itemize}

Предположение: с увеличением количества узлов, время выполнения анализа будет расти.

\subsubsection{Результаты исследования}
В таблице~\ref{tbl:exp-1} представлена выжимка данных, полученных в результате проведения первого исследования.
Полные данные, полученные в результате проведения исследования, представлены в таблице~\ref{tbl:exp-1-full}. 
На рисунке~\ref{img:cascade.pdf} представлен график зависимости времени выполнения анализа
от количества узлов системы.
\begin{center}
    \begin{tabularx}{\textwidth}{|X|X|}
        \caption{\label{tbl:exp-1}Результаты проведения первого исследования} \\
        \hline
        Количество узлов & Время выполнения анализа, мс \\ \hline
        \endfirsthead

        \multicolumn{2}{l}{{Продолжение таблицы\ \thetable{}}} \\ \hline
        Количество узлов & Время выполнения анализа, мс \\ \hline
        \endhead

        \hline
        \endfoot

        \hline
        \endlastfoot

        1000 & 3.5 \\ \hline
        2000 & 3.7 \\ \hline
        3000 & 6 \\ \hline
        4000 & 7.5 \\ \hline
        5000 & 16.4 \\ \hline
        6000 & 7.6 \\ \hline
        7000 & 11.9 \\ \hline
        8000 & 15.4 \\ \hline
        9000 & 14.5 \\ \hline
        10000 & 15.5 \\ \hline
    \end{tabularx}
\end{center}
\jimg{cascade.pdf}{Зависимость времени выполнения анализа от количества узлов системы}

\subsubsection{Вывод}
Несмотря на зашумленность графика, прослеживается зависимость времени выполнения
анализа от количества узлов системы: время выполнения незначительно увеличивается при росте
количества узлов: при десятикратном увеличении количества узлов, время выполнения анализа
выросло в 3 раза. Абсолютные значения времени достаточно малы ---
единицы и десятки миллисекунд, что позволяет использовать разработанную систему и БД
на графах с тысячами вершин. 

\subsection{Второе исследование}
\subsubsection{Описание исследования}
Данное исследование направлено на выявление зависимости времени выполнения поиска
циклических зависимостей от количества циклов в системе. Условия исследования:
\begin{itemize}
    \item количество циклов в графе изменяется от 10 до 50 с шагом 1;
    \item количество узлов рассчитывается как \( 50 + C_{cnt} \cdot 3 \), где \( C_{cnt} \) --- количество циклов в графе;
    \item количество связей в графе системы равняется количеству узлов, умноженному на два, на каждом шаге.
\end{itemize}

Предположение: с увеличением количества циклов, время выполнения анализа будет расти.

\subsubsection{Результаты исследования}
В таблице~\ref{tbl:exp-2} представлена выжимка данных, полученных в результате 
проведения второго исследования. Полные данные, полученные в результате 
проведения исследования, представлены в таблице~\ref{tbl:exp-2-full}. 
На рисунке~\ref{img:cycles.pdf} представлен график зависимости 
времени подсчета циклов от количества циклов в графе.
\begin{center}
    \begin{tabularx}{\textwidth}{|X|X|}
        \caption{\label{tbl:exp-2}Результаты проведения второго исследования} \\
        \hline
        Количество циклов & Время подсчета циклов, мс \\ \hline
        \endfirsthead

        \multicolumn{2}{l}{{Продолжение таблицы\ \thetable{}}} \\ \hline
        Количество циклов & Время подсчета циклов, мс \\ \hline
        \endhead

        \hline
        \endfoot

        \hline
        \endlastfoot

        10 & 8.4 \\ \hline
        15 & 10.7 \\ \hline
        20 & 13.0 \\ \hline
        25 & 13.9 \\ \hline
        30 & 22.2 \\ \hline
        35 & 26.2 \\ \hline
        40 & 32.2 \\ \hline
        45 & 38.5 \\ \hline
        50 & 41.5 \\ \hline
    \end{tabularx}
\end{center}

\clearpage

\jimg{cycles.pdf}{Зависимость времени подсчета циклов от количества циклов в графе}

\subsubsection{Вывод}
На представленном графике прослеживается зависимость времени поиска циклов 
в графе от количества циклов в системе: время выполнения 
увеличивается при росте количества циклов: при десятикратном увеличении 
количества узлов, время выполнения анализа выросло в 5 раз. Абсолютные 
значения времени достаточно малы --- и десятки миллисекунд, что 
позволяет использовать разработанную систему и БД на графах с сотнями вершин
и десятками циклов. 

\numberlesssubsection{Выводы по исследовательскому разделу}
В данном разделе была проведена оценка производительности 
выбранной графовой базы данных при работе с топологиями различной сложности. 
Операции глубокого обхода дерева для выявления каскадных зависимостей и поиск 
всех замкнутых контуров в ориентированном графе относятся к классу вычислительно 
сложных задач. В традиционных системах хранения такие алгоритмы требуют значительных 
вычислительных мощностей и времени на выполнение рекурсивных проверок.

Результаты проведенных исследований продемонстрировали высокую временную 
эффективность используемой графовой системы управления базами данных. При 
масштабировании инфраструктуры до десяти тысяч узлов, а также при значительном 
увеличении количества циклических связей в системе, время обработки аналитических 
запросов возрастает предсказуемо и остается в пределах нескольких десятков миллисекунд.